\labsection{5}{Leap years, cubes, and Armstrong numbers}{sec:lab05}

\begin{labproject}{Lab 5b guided practice}
\textbf{Coverage.} Chapter~\ref{chap:control-flow}.\\
\textbf{Goal.} Solve three problems with conditions and loops.

\textbf{Task 1 --- Check whether a year is a leap year.}

\begin{lstlisting}[style=dsr]
year <- as.integer(readline("Input year: "))

if ((year %% 400 == 0) ||
    (year %% 4 == 0 && year %% 100 != 0)) {
  print(paste(year, "is a leap year."))
} else {
  print(paste(year, "is not a leap year."))
}
\end{lstlisting}

The sample classifies 2004 as a leap year and 1900 as a common year.

\textbf{Task 2 --- Display cubes up to an input integer.}

\begin{lstlisting}[style=dsr]
n <- as.integer(readline("Input an integer: "))

for (i in seq_len(n)) {
  cube <- i ^ 3
  print(paste(
    "Number is:", i,
    "and cube of the", i, "is:", cube
  ))
}
\end{lstlisting}

For \texttt{n = 5}, the cubes are 1, 8, 27, 64, and 125.

\textbf{Task 3 --- Check an Armstrong number.}

An Armstrong number equals the sum of its digits raised to the number of digits. This definition correctly classifies the supplied example \texttt{1634}.

\begin{lstlisting}[style=dsr]
number <- as.integer(readline("Input an integer: "))

original <- number
digits <- nchar(as.character(number))
sum_power <- 0

temp <- number
while (temp > 0) {
  digit <- temp %% 10
  sum_power <- sum_power + digit ^ digits
  temp <- temp %/% 10
}

if (sum_power == original) {
  print(paste(original, "is an Armstrong number."))
} else {
  print(paste(original, "is not an Armstrong number."))
}
\end{lstlisting}

\begin{pitfall}
The source also says ``power three,'' which conflicts with its four-digit example. Confirm the required definition with your instructor.
\end{pitfall}
\end{labproject}

\begin{labdeliverable}
Submit all three programs with representative test results.
\end{labdeliverable}
